\section{Expositor of Scripture, Preacher in the Ruins, and Author of the \work{Dialogi}}

\subsection{\work{Moralia in Iob}: the three senses and the stricken expositor}

Alongside the \work{Regula Pastoralis}, Gregory's greatest doctrinal and
spiritual legacy lies in his exposition of Holy Scripture. His vast
\work{Moralia in Iob}---thirty-five books arranged in six tomes, begun as oral
conferences to his monks and Leander in Constantinople, continued by dictation,
and revised for publication around 595 (\work{Ep.}\ I.43, V.49)---became one
of the foundational textbooks of medieval theology, asceticism, and moral
psychology.

In the Preface and Epistle Dedicatory to Leander, Gregory lays out a serene,
organic theology of Scripture. Asking who humanly penned the Book of Job, he
answers, is like receiving a letter from a great man and fussing over which
pen wrote the ink: ``When then we understand the matter, and are persuaded that
the Holy Spirit was its Author, in stirring a question about the author, what
else do we than in reading a letter enquire about the pen?'' (\work{Moralia},
Preface 2). He approaches the sacred page with a disciplined threefold method
that shaped Western biblical interpretation for centuries:

\begin{quote}
For first, we lay the historical foundations; next, by pursuing the typical
sense, we erect a fabric of the mind to be a stronghold of faith; and moreover
as the last step, by the grace of moral instruction, we, as it were, clothe the
edifice with an overcast of colouring. (\work{Moralia in Iob}, Ep.\ Ded.\ 3)
\end{quote}

Gregory was neither a rigid literalist nor a reckless allegorizer. Some verses
(like Job's lament cursing the day of his birth) cannot stand at a crude
surface level and beckon the reader deeper; yet others---such as Job's care
for the widow, the orphan, and the freezing poor (Job 31:16--20)---must never
be dissolved into mere allegory, lest ``we make void all his acts of mercy''
(\work{Moralia}, Ep.\ Ded.\ 4). Scripture accommodates every soul at once:

\begin{quote}
It is, as it were, a kind of river, if I may so liken it, which is both
shallow and deep, wherein both the lamb may find a footing, and the elephant
float at large. (\work{Moralia in Iob}, Ep.\ Ded.\ 4)
\end{quote}

The interpreter, too, should follow ``the way of a river'' that turns aside
into open side valleys whenever an opportunity for moral edification appears,
before flowing back into its main channel (Ep.\ Ded.\ 2). Gregory openly
commented on Jerome's ``new Translation'' (the Vulgate) while freely citing the
Old Latin whenever helpful, because ``the Apostolic See, over which I preside
by ordinance of God, uses both'' (Ep.\ Ded.\ 5). And he interpreted Job not
from comfortable academic leisure, but from the furnace of chronic pain:
``perchance it was this that Divine Providence designed, that I a stricken one,
should set forth Job stricken, and that by these scourges I should the more
perfectly enter into the feelings of one that was scourged'' (Ep.\ Ded.\ 5).
Through thirty-five books, Gregory wove Augustine's theology of grace and
interiority into a contemplative school of compunction---the piercing sorrow
of love that turns a bruised soul away from passing noise toward the face of
God.

\subsection{\work{Homiliae XL in Evangelia}: preaching across the stational basilicas}

Between late 590 and 592, Gregory delivered forty Gospel homilies ``amid the
sacred solemnities of the Mass'' (\latin{inter sacra missarum solemnia}) across
Rome's stational basilicas---Saint Peter's, Saint Mary Major, Saint Clement's,
Saint Andrew's, Saint Stephen's, Saint Felicitas's, and the Lateran. In his
preface sending the corrected edition to Bishop Secundinus of Taormina in
Sicily, Gregory describes the exact mechanics of their composition: the first
twenty homilies he dictated beforehand and had read aloud by a notary while he
sat present (often incapacitated by illness); the second twenty he spoke
directly before the congregation while stenographers took down his words.
Certain eager brothers, ``like famished men who try to eat food before it is
fully cooked,'' had begun copying his drafts before he could revise them, so
Gregory deposited the authorized two-volume text in the papal archive
(\latin{in scrinio sanctae Ecclesiae nostrae}) as the standard (\work{Homiliae
in Evangelia}, Praefatio).

The homilies show Gregory speaking directly to ordinary Roman townspeople with
extraordinary concreteness, tenderness, and eschatological urgency. On the
Second Sunday of Advent in Saint Peter's Basilica (\work{Homiliae in
Evangelia} 1.1, 1.5), preaching on Luke 21:25--32, he looked out over a
shrinking congregation battered by plague, storms, and war and drew a haunting
portrait of an aging world (\latin{senescentem mundum}):

\begin{quote}
Behold, my brothers, we now see what we used to hear. Every day the world is
pressed by new and multiplying evils. You see how few of you remain out of that
countless populace; and yet scourges still press upon us daily, sudden blows
strike us down, new and unforeseen disasters afflict us. \dots\ So the world in
its earlier years flourished as in youth, strong to bring forth the human race,
fresh in bodily health, rich in abundance of goods; but now it is bowed down by
its very old age (\latin{ipsa sua senectute deprimitur}), and pressed by
growing troubles as toward a nearing death. Do not, therefore, my brothers,
love that which you see cannot stand long. (\work{Homiliae in Evangelia} 1.5)
\end{quote}

Yet Gregory's preaching was never mere gloom; it glowed with sacramental
wonder. Preaching on Christmas Day (\work{Homiliae in Evangelia} 8.1), he
opened with a precious liturgical notice---``Because, by the Lord's bounty, we
are to celebrate the solemnities of Mass three times today, we cannot speak
long on the Gospel reading''---before fixing his hearers' eyes on one piercing
etymology:

\begin{quote}
Rightly also is He born in Bethlehem: for Bethlehem is interpreted ``the house
of bread'' (\latin{domus panis}). He Himself is the one who says, ``I am the
living bread which came down from heaven.'' The place where the Lord is born
was therefore called beforehand the house of bread, because it was surely to
come to pass that He would appear there in the substance of flesh who would
refresh the minds of the elect with inward fullness. \dots\ And because it is
said through the prophet, ``All flesh is grass,'' being made man He turns our
grass into wheat (\latin{fenum nostrum vertit in frumentum}), who says of
Himself, ``Unless a grain of wheat falling into the earth die, it remains
alone.'' Therefore, newly born, He is laid in a manger, that He might refresh
all the faithful---the holy beasts---with the wheat of His flesh, lest they
remain fasting from the nourishment of eternal understanding. (\work{Homiliae
in Evangelia} 8.1)
\end{quote}

At the close of many homilies, Gregory added vivid exempla from contemporary
Roman memory to turn doctrine into action: his aunts Tharsilla, Aemiliana, and
Gordiana (38.15); the poor paralyzed virgin Romula living beside Saint Mary
Major, ``a pearl of God hidden on a dunghill'' whose passing filled her cell
with light and heavenly psalmody (40.11--12); and the captive husband whose
iron chains fell loose each week at the very hour his wife offered the
Eucharistic sacrifice for him, thinking him dead (37.8). In that homily Gregory
articulated his deep Eucharistic realism: ``as often as we offer to Him the
victim of His passion, so often we renew His passion for our absolution''
(\work{Homiliae in Evangelia} 37.7).

\subsection{\work{Homiliae in Hiezechihelem}: preaching beneath Lombard swords (593)}

In 593, Gregory turned from the Gospels to the prophet Ezekiel, preaching
twenty-two homilies in two books (twelve on the vision of the chariot and the
prophet's call; ten on the mystical temple of Ezekiel 40). Here Gregory's
contemplative theology reaches its deepest maturity---explaining how the
contemplative soul sees more than language can carry (II.6.1), how faith is
the threshold leading upward through the seven gifts of the Spirit (II.7.7--9),
and how Martha and Mary must coexist inside every Christian heart.

Yet history broke violently into the nave. As Lombard King Agilulf marched
down from Pavia toward Rome in late 593, Gregory stood in the pulpit gazing on
a dying imperial capital. In Book II, Homily 6, commenting on Ezekiel's
boiling cauldron emptied over the coals, he uttered the most unforgettable
elegy over the fall of ancient Rome in all of Latin literature:

\begin{quote}
Where is the senate? Where now is the people? (\latin{Ubi enim senatus? Ubi
jam populus?}) The bones have wasted away, the flesh is consumed, all the pride
of worldly dignities in her is extinguished. \dots\ Because the senate is gone,
the people has perished, and yet among the few who remain sorrows and groans
multiply every day, Rome now burns empty (\latin{jam vacua ardet Roma}). Why do
we speak of men, when with multiplying ruins we see the very buildings
destroyed? \dots\ For after the men have failed, even the walls are falling
(\latin{postquam defecerunt homines, etiam parietes cadunt}). (\work{Homiliae
in Hiezechihelem} II.6.22)
\end{quote}

Four homilies later, with Agilulf's army outside the gates, Gregory could go
no further. The closing lines of Book II, Homily 10 record the exact moment
the siege silenced the preacher:

\begin{quote}
Let no one blame me if after this discourse I cease, because, as you all see,
our tribulations have grown beyond measure: on every side we are surrounded by
swords (\latin{undique gladiis circumfusi sumus}), on every side we fear the
imminent peril of death. Some return to us with their hands cut off; others are
reported captured, others slain. Now I am forced to withhold my tongue from
exposition, for ``my soul is weary of my life'' (Job 10:1). \dots\ For him who
is not permitted to live, how can it please him to speak of the mystical senses
of sacred Scripture? And I who am forced every day to drink bitter draughts,
when can I pour out sweet ones? What then remains, save that amid the scourges
which we suffer for our iniquities we give thanks with tears? (\work{Homiliae
in Hiezechihelem} II.10.24)
\end{quote}

\subsection{The \work{Dialogi}: Italian sanctity and the world beyond}

Between 593 and 594, in the immediate aftermath of the siege, Gregory composed
the four books of the \work{Dialogi}---conversations with his deacon Peter on
the lives and miracles of the holy men of Italy (summarized in Barmby's
Prolegomena and cross-referenced in Migne PL 76).

Book II of the \work{Dialogi} is the sole early biography of Saint Benedict of
Nursia, whose \work{Regula Benedicti} Gregory knew and revered, transmitting
Benedict's memory from Subiaco and Monte Cassino across Western monasticism.
Book IV addresses death, the survival of the soul, purifying fire after death
for lighter faults, and the power of the Eucharistic sacrifice offered for the
departed---including the celebrated case of the monk Justus of Saint Andrew's
monastery (\work{Dialogi} IV.55/57).

Modern readers trained on the sober administrative prose of the
\work{Registrum} have sometimes been startled by the abundance of visions and
wonders in the \work{Dialogi}, and in 1987 Francis Clark famously argued that
the work was a late seventh-century forgery. Patristic scholarship, however,
has overwhelmingly rejected that hypothesis and reaffirmed Gregory's
authorship. Gregory's own authenticated writings confirm why: the very same
miracle narratives---his aunts Tharsilla and Aemiliana, the virgin Romula,
Bishop Cassius of Narni, and the captive whose fetters broke during
Mass---appear almost verbatim inside his \work{Homiliae in Evangelia}
(37.8--9, 38.15--16, 40.11), and his letter to Empress Constantina (\work{Ep.}\
IV.30) expresses the exact same vivid conviction about terrifying prodigies
guarding the tombs of Saints Peter, Paul, and Lawrence and the miraculous power
of contact relics (\latin{brandea}) and filings from Saint Peter's chains. For
Gregory, recording Italian miracles was pastoral consolation: proving to a
terrified, Lombard-ravaged people that God had not abandoned Italy, and that
holiness still walked its hills.
